// Wood Grain — Concentric rings with FBM perturbation
//
// Generates realistic wood-like rings by computing distance from a
// centre point, multiplying by a ring frequency, then perturbing the
// result with FBM noise. A colour ramp maps from light to dark wood.
//
// Inputs:  (none — purely procedural)
// Outputs: @OUT (vec3)

f$ring_freq   = 20.0;    // number of rings across the surface
f$noise_scale = 4.0;     // scale of the FBM perturbation
f$noise_amp   = 2.5;     // strength of the FBM perturbation
f$center_x    = 0.5;     // ring centre X (0-1 UV space)
f$center_y    = 0.5;     // ring centre Y (0-1 UV space)

float x = u * $noise_scale;
float y = v * $noise_scale;

// Distance from ring centre
float dist = distance(vec2(u, v), vec2($center_x, $center_y));

// Perturb with FBM for organic variation
float perturbation = fbm(x, y, 4) * $noise_amp;

// Concentric rings
float ring = sin(dist * $ring_freq + perturbation);

// Remap sin output from [-1,1] to [0,1]
float t = ring * 0.5 + 0.5;

// Wood colour palette: light sapwood to dark heartwood
vec3 light_wood = vec3(0.85, 0.65, 0.35);
vec3 mid_wood   = vec3(0.55, 0.35, 0.15);
vec3 dark_wood  = vec3(0.30, 0.18, 0.08);

// Two-segment ramp
vec3 result = t < 0.5
    ? lerp(dark_wood, mid_wood, t * 2.0)
    : lerp(mid_wood, light_wood, (t - 0.5) * 2.0);

@OUT = result;
